\section{Introdução}
\label{sc:introducao}

Nos últimos anos, mesmo com a queda nos índices de desmatamento no Brasil, a ocorrência de queimadas não diminuiu — impulsionados por secas severas, os incêndios atingiram níveis recordes na Amazônia, Cerrado e Pantanal. O avanço do fogo afeta diversas frentes: no âmbito socioambiental, mais de 160 Terras Indígenas registraram focos de calor, incluindo o incêndio que destruiu 38 moradias na Aldeia Canoanã, na Ilha do Bananal (TO) \cite{g1_canoana_2026}; na saúde e logística, a fumaça densa elevou atendimentos respiratórios no Sistema Único de Saúde (SUS) e interditou rodovias federais; e no setor elétrico, o calor e a fuligem sob linhas de transmissão do Sistema Interligado Nacional (SIN) reduzem o isolamento do ar, causando curtos-circuitos (\textit{flashover}) e quedas de energia.

Embora o Instituto Nacional de Pesquisas Espaciais (INPE) forneça dados abertos de satélites \cite{inpe2025bdqueimadas}, cada registro bruto resume-se a uma coordenada isolada (latitude e longitude) que não informa qual comunidade ou infraestrutura está em risco. Além disso, cruzar milhões de coordenadas contra mapas vetoriais no momento da consulta torna bancos relacionais lentos. Para solucionar esse problema, desenvolveu-se na Iniciação Científica (IC) a plataforma \textbf{SIRDA 1.0} \cite{sampaio2026ic}: uma arquitetura que antecipa o cruzamento espacial em memória e agrupa focos vizinhos em perímetros de incêndio (via ST-DBSCAN) antes de transmitir os eventos via Apache Kafka para o banco colunar Apache Druid \cite{yang2014druid} e painéis no Apache Superset. Ao todo, unificaram-se 4,05 milhões de registros de 7 bases públicas (INPE, ONS, ANEEL, FUNAI, CNES/DATASUS, DNIT e IBGE), monitorando hospitais do SUS, terras indígenas, rodovias federais, municípios e a rede elétrica por bioma.

A operação da plataforma sobre um ano completo de dados validou a estabilidade da arquitetura e revelou padrões sazonais de pressão térmica com consultas em frações de segundo \cite{sampaio2026ic}. A partir dessa base consolidada, este Trabalho de Conclusão de Curso (TCC) propõe evoluir o sistema para o \textbf{SIRDA 2.0}, integrando modelos de linguagem (via protocolo FastMCP e busca documental RAG), automação de fluxos analíticos e alertas (Webhooks), processamento analítico híbrido (DuckDB \cite{raasveldt2019duckdb} e Apache Druid) e modelagem direcional de propagação do fogo para apoiar a tomada de decisão rápida baseada em dados em infraestruturas estratégicas.
